% Successor v10 M2 high-source canonical deletion guard (2026-10-12).
% New independent normalized Lean endpoint PENDING full/focused exact-head CI.
% No mathematical author prose or four prior corrections are modified.
%
\todo[inline]{\textbf{v10 M2: safe canonical deletion above the boundary
-- PENDING exact-head Lean verification.}
Let $F$ be shape-preserving with a strict gap before the first
affected source level $n$, let $m=\widetilde F(n)-1$, and consider
a genuine source successor at level $j\geq n$.
Write its mapped canonical successor predecessor level as
$k+1=\widetilde F(j)$, and its canonical new E cut as $f$.
The proposed theorem
\texttt{mapped\_upper\_canonical\_step\_safe\_for\_gap}
shows \emph{either} $m<k$ \emph{or} $f<m$.
If $j>n$, strict preservation of levels makes the first
alternative hold. At $j=n$, (S1), the image target gap and
the actual empty-or-singleton canonical parameter record
force the second: $f=0$ when the parameter is empty;
otherwise the unique parameter has level $f<m$.
The old parameter list is exactly the image of the genuine
source successor parameters, without an additional age,
successor or free-cut hypothesis.

This proves the exact condition needed for
\texttt{canonicalKptStep\_of\_deleted\_record}
to reconstruct a genuine successor after deleting $m$,
including the boundary step where the preceding ordinary
vertex would otherwise become consecutive in $E$.
It does \emph{not} itself prove weak successor preservation
of the intrinsic deleted-image function, injection or
M2 factorization; corrected finite $L$-only B3 necessity
remains OPEN.}
